\section{Pattern Interpretability and Actionability}
\label{sec:2_7}

Discovered patterns have value only if they can be understood and operationalized. This section reviews methods for interpreting patterns and translating them into investigation guidance.

\subsection{Explaining Discovered Patterns}

\textbf{SHAP} (SHapley Additive exPlanations) by \citeauthor{lundberg2017unified} \cite{lundberg2017unified} provides theoretically grounded feature attribution through Shapley values from cooperative game theory. For a prediction, SHAP computes each feature's contribution satisfying local accuracy, consistency, and missingness properties. TreeSHAP \cite{lundberg2020local} enables efficient computation for XGBoost, directly applicable to the validation classifier in hybrid approaches.

SHAP produces complementary explanation types:

\textbf{Global explanations} aggregate feature importance across the dataset, revealing which patterns most strongly predict fraud. Beeswarm plots show each feature's impact distribution, supporting model auditing and pattern validation.

\textbf{Local explanations} decompose individual predictions into feature contributions. For a flagged listing, analysts see exactly which patterns elevated fraud score ``shared device with known fraud: +15\%, high-risk IP country: +12\%'' enabling informed investigation prioritization.

\textbf{GNN explanation methods} address the challenge of interpreting relational patterns encoded in embeddings. GNNExplainer \cite{ying2019gnnexplainer} identifies subgraph structures most influential for predictions through mutual information maximization. PGExplainer \cite{luo2020parameterized} parameterizes explanation generation for collective pattern identification across instances. GraphSVX \cite{duval2021graphsvx} extends Shapley values to graphs, aligning with SHAP methodology used for the validation classifier.

\subsection{Operational Requirements}

Beyond technical capability, explanations must meet operational requirements.

\textbf{Analyst trust} develops when explanations align with domain expertise. If patterns match known fraud indicators (suspicious countries, device sharing, disposable emails), trust builds. Unexpected patterns prompt investigation revealing either novel fraud tactics or model deficiencies.

\textbf{Actionability} requires explanations that guide investigation. Knowing ``GNN dimension 7 contributed +8\%'' provides little guidance; knowing ``shared device with 3 flagged listings: +8\%'' directs analyst attention to specific relationships.

\textbf{Consistency} between global and local explanations supports system reliability. If global analysis identifies device sharing as the top pattern, individual explanations should generally reflect this factor.

This thesis integrates SHAP for the XGBoost classifier, quantifying contributions of both individual attributes and aggregated relational pattern representations (composite features).
